\documentclass[11pt]{article}
\usepackage[a4paper,margin=1in]{geometry}
\usepackage{amsmath,amssymb,amsthm,booktabs,mathtools}
\usepackage[T1]{fontenc}
\usepackage{lmodern}
\usepackage{hyperref}
\setlength{\emergencystretch}{3em}

\newtheorem{theorem}{Theorem}[section]
\newtheorem{proposition}[theorem]{Proposition}
\newtheorem{corollary}[theorem]{Corollary}
\newtheorem{remark}[theorem]{Remark}
\newcommand{\Mob}{\mu}
\newcommand{\SigmaA}{\Sigma_A}
\newcommand{\Cov}{\operatorname{Cov}}
\newcommand{\Var}{\operatorname{Var}}
\newcommand{\E}{\mathbb E}
\newcommand{\R}{\mathbb R}

\title{A branch-symmetric Markov family on the RH-366 graph:\\
exact Möbius orthogonality and covariance}
\author{RH research program}
\date{August 2026}

\begin{document}
\maketitle

\begin{abstract}
We derive, rather than geometrically select, a one-parameter Markov/Gibbs
family on the primitive four-state graph used by RH-366.  In the state order
$(--,-+,+-,++)$, put $q=1-t$ and
\[
 P_t=\begin{pmatrix}
 t&0&q&0\\
 1&0&0&0\\
 0&t&0&q\\
 0&1&0&0
 \end{pmatrix},\qquad 0<t<1.
\]
Its stationary law, centered sign variance, and all-lag covariance are exact:
the odd covariances vanish and the variance-one covariance at lag $2k$ is
$(-q)^k$.  Each fixed parameter gives a mixing Markov measure, so the
finite-state Chebyshev--Borel--Cantelli argument yields Möbius orthogonality
almost surely, simultaneously for all continuous observables.  The weighted
finite variance has an exact even-shift expansion and a geometric bound.
The Parry law is the single member $t=\varphi^{-1}$; the open interval adds
non-Parry measures but no canonical arithmetic coupling.  All statements are
fixed-parameter, and no uniform endpoint theorem, operator trace, zero model,
or RH implication is claimed.
\end{abstract}

\section{Frozen graph and derived data type}

The source package
\cite{HenonMobius2026} freezes the state order
\[
 (--,-+,+-,++)
\]
and the adjacency matrix
\[
 A=\begin{pmatrix}
 1&0&1&0\\
 1&0&0&0\\
 0&1&0&1\\
 0&1&0&0
 \end{pmatrix}.
 \tag{1}
\]
The first coordinate of a state is the scalar sign observable
\[
 e=(-1,-1,+1,+1)^{\mathsf T}. \tag{2}
\]
The graph and its Hénon conjugacy are inherited inputs; they are not
reconstructed from the finite calculations below.

For $0<t<1$, let $q=1-t$ and define the row-stochastic matrix
\[
 P_t=\begin{pmatrix}
 t&0&q&0\\
 1&0&0&0\\
 0&t&0&q\\
 0&1&0&0
 \end{pmatrix}. \tag{3}
\]
Every positive entry of $P_t$ lies on an edge of (1), and every edge of
(1) receives a positive probability.  Thus (3) is a new derived symbolic
family, not a source claim that the Hénon geometry chooses a parameter.

\begin{proposition}[Stationarity and mixing]
\label{prop:mixing}
Put $d=1+q=2-t$.  The unique stationary row vector of $P_t$ is
\[
 \pi_t=\frac{(1,q,q,q^2)}{d^2}. \tag{4}
\]
Moreover,
\[
 \det(\lambda I-P_t)=(\lambda-1)(\lambda+q)(\lambda^2+q). \tag{5}
\]
Consequently $P_t$ is primitive and its nontrivial spectral radius is
$\sqrt q<1$.
\end{proposition}

\begin{proof}
Multiplying the row in (4) by (3) gives the same row, and its entries sum to
one.  The directed support is the primitive graph (1); in fact a direct
multiplication gives $A^4>0$, and all the corresponding entries of $P_t^4$
are positive for $0<t<1$.  This proves primitivity and uniqueness of the
stationary law.  Expanding the determinant in (5) gives the stated
factorization.  The eigenvalue $1$ is simple, while the other three
eigenvalues are $-q$ and $\pm i\sqrt q$.
\end{proof}

\section{Exact observable and covariance}

Let $\nu_t$ be the stationary Markov measure on $\SigmaA$ associated with
$P_t$.  With $e$ from (2), direct multiplication gives
\[
 m_t=\pi_t e=-\frac{t}{d},\qquad
 v_t=\pi_t(e-m_t{\bf 1})^2=\frac{4q}{d^2}. \tag{6}
\]
Since $q>0$, define the centered variance-one observable
\[
 F_t=\frac{e-m_t{\bf 1}}{\sqrt{v_t}}
 =(-\sqrt q,-\sqrt q,q^{-1/2},q^{-1/2})^{\mathsf T}. \tag{7}
\]

\begin{proposition}[All-lag covariance]
\label{prop:cov}
For every $k\geq0$,
\[
 \Cov_{\nu_t}(F_t,F_t\circ\sigma^{2k+1})=0,\qquad
 \Cov_{\nu_t}(F_t,F_t\circ\sigma^{2k})=(-q)^k. \tag{8}
\]
For the raw sign, the even covariance is $v_t(-q)^k$ and the odd
covariance is zero.
\end{proposition}

\begin{proof}
The stationary covariance formula for a state observable is
\[
 \Cov_{\nu_t}(F_t,F_t\circ\sigma^\ell)
 =\pi_t\operatorname{diag}(F_t)P_t^\ell F_t. \tag{9}
\]
Using (3), (4), and (7), one obtains the two exact identities
\[
 P_t^2F_t=-qF_t,\qquad
 \pi_t\operatorname{diag}(F_t)P_tF_t=0. \tag{10}
\]
For even $\ell=2k$, (9)--(10) give $(-q)^k$.  For odd
$\ell=2k+1$, factor out $(-q)^k$ and use the second identity in (10).
\end{proof}

\begin{remark}[The Parry point]
For $\varphi=(1+\sqrt5)/2$, the RH-366 Parry matrix is exactly (3) at
\[
 t=\varphi^{-1},\qquad q=\varphi^{-2}. \tag{11}
\]
Thus (8) recovers the RH-366 covariance.  The theorem is new only in the
parameterized non-Parry family and in the fixed-parameter quantifier; it does
not rebrand the Parry calculation as a physical spectral result.
\end{remark}

\section{Fixed-parameter Möbius orthogonality}

\begin{theorem}[Almost-sure orthogonality for each fixed parameter]
\label{thm:as}
For every fixed $t\in(0,1)$ there is a set $E_t\subset\SigmaA$ with
$\nu_t(E_t)=1$ such that, for every $\omega\in E_t$ and every
$f\in C(\SigmaA)$,
\[
 \frac1N\sum_{n\leq N}\Mob(n)f(\sigma^n\omega)\longrightarrow0. \tag{12}
\]
Under the frozen Hénon conjugacy, the same statement holds for continuous
observables on the local survivor.
\end{theorem}

\begin{proof}
It is enough first to treat a centered locally constant function $g$.  The
primitive finite-state chain has exponential mixing, so there are constants
$C_g<\infty$ and $0<\rho_t<1$ such that
\[
 \left|\Cov_{\nu_t}\bigl(g\circ\sigma^i,g\circ\sigma^j\bigr)\right|
 \le C_g\rho_t^{|i-j|}. \tag{13}
\]
With $T_N(\omega)=\sum_{n\le N}\Mob(n)g(\sigma^n\omega)$ and
$|\Mob(n)|\le1$, (13) implies
\[
 \E_{\nu_t}|T_N|^2\le C'_g N. \tag{14}
\]
Chebyshev's inequality at $N=k^2$ gives a summable bound for
\(\nu_t\{|T_{k^2}|>\epsilon k^2\}\).  Borel--Cantelli yields
$T_{k^2}/k^2\to0$ almost surely.  Between consecutive squares there are
$O(k)$ bounded summands, so the convergence holds for all $N$.

Take a countable dense family of centered cylinder functions with rational
values and intersect the corresponding full-measure sets.  Uniform
approximation extends the conclusion to all centered continuous functions.
For a general $f$, subtract its $\nu_t$-mean.  The remaining constant term is
a multiple of $\sum_{n\le N}\Mob(n)=o(N)$, by the prime number theorem.
\end{proof}

\begin{remark}[Quantifier firewall]
The set $E_t$ and the constants in (13)--(14) depend on the fixed parameter.
No common full-measure set over uncountably many $t$, and no estimate uniform
as $t\downarrow0$ or $t\uparrow1$, is asserted.
\end{remark}

\section{Möbius-weighted variance}

Define
\[
 S_{N,t}(\omega)=\sum_{n\le N}\Mob(n)F_t(\sigma^n\omega),
 \qquad V_{N,t}=\Var_{\nu_t}(S_{N,t}). \tag{15}
\]
The covariance table gives the exact identity
\[
 V_{N,t}
 =\sum_{n\le N}\Mob(n)^2
 +2\sum_{k=1}^{\lfloor(N-1)/2\rfloor}(-q)^k
   \sum_{n\le N-2k}\Mob(n)\Mob(n+2k). \tag{16}
\]

\begin{proposition}[Unconditional finite bound]
\label{prop:bound}
For every $N\ge1$ and $0<t<1$,
\[
 0\le V_{N,t}\le \left(1+2\sum_{k\ge1}q^k\right)N
 =\frac{2-t}{t}\,N. \tag{17}
\]
Equivalently,
\[
 \left\|\frac{S_{N,t}}N\right\|_{L^2(\nu_t)}
 \le\sqrt{\frac{2-t}{tN}}. \tag{18}
\]
\end{proposition}

\begin{proof}
Nonnegativity is the variance definition.  In (16), bound the diagonal by
$N$, each inner correlation by $N$, and sum the geometric absolute values.
The $L^2$ estimate is the square root of (17) divided by $N$.
\end{proof}

\begin{theorem}[Conditional variance density]
\label{thm:conditional}
Assume the ordinary fixed-shift two-point correlations
\[
 \frac1N\sum_{n\le N}\Mob(n)\Mob(n+2k)\longrightarrow0
 \quad\text{for every fixed }k\ge1. \tag{19}
\]
Then, for every fixed $t\in(0,1)$,
\[
 \frac{V_{N,t}}N\longrightarrow\frac6{\pi^2}. \tag{20}
\]
\end{theorem}

\begin{proof}
The squarefree density gives $N^{-1}\sum_{n\le N}\Mob(n)^2\to6/\pi^2$.
For the off-diagonal part of (16), first retain finitely many $k$ and use
(19).  The remaining tail is bounded uniformly in $N$ by
$2\sum_{k>K}q^k$, which tends to zero as $K\to\infty$.  This proves (20).
\end{proof}

\begin{remark}[Scope of the conditional statement]
Equation (20) is not an unconditional two-point Chowla theorem.  It is a
conditional consequence recorded to identify exactly which arithmetic input
would be needed for the variance density.  The unconditional theorem is
(16)--(18).
\end{remark}

\section{Parameter boundaries and route decision}

At $t=0$, the branch choices disappear and the chain is a deterministic
period-four system; at $t=1$, the stationary law collapses and $v_t=0$, so
$F_t$ is undefined.  Also, the nontrivial spectral radius
$\sqrt{1-t}$ tends to one as $t\downarrow0$, while the entries
$q^{-1/2}$ in (7) diverge as $t\uparrow1$.  These are genuine reasons to keep
the theorem pointwise on the open interval.

The stationary entropy rate, included only as a structural diagnostic, is
\[
 h(t)=\frac{-t\log t-(1-t)\log(1-t)}{2-t}. \tag{21}
\]
Its maximum occurs at the Parry point (11); this does not select a canonical
parameter for the arithmetic problem.

\begin{center}
\begin{tabular}{p{0.23\linewidth}p{0.68\linewidth}}
\toprule
Route & Verdict and boundary\\
\midrule
Route A & \texttt{GO}.  The derived family has exact stationary, mixing,
covariance, almost-sure fixed-parameter Möbius, and finite-variance theorems.\\
Route B & \texttt{STOP\_SCOPED}.  The parameter is externally chosen; no
canonical arithmetic coupling, determinant, prime-power trace, zeta-zero
model, or physical Gate is supplied.\\
\bottomrule
\end{tabular}
\end{center}

No Hilbert--Pólya operator is constructed, no Riemann zero is identified, and
the Riemann hypothesis is not addressed.

\section{Executable audit and overlap ledger}

The package checks stationarity, the centered $P_t^2$ eigenidentity, the
vanishing one-step covariance, positivity of $P_t^4$, and the characteristic
covariance table at rational parameters $1/2$, $2/3$, and $3/4$.  It also
compares (16) with an independent finite double sum through $N=14$.
These are exact finite checks, not asymptotic fits.

The overlap is narrow.  RH-366 supplies the graph, sign observable, Parry
specialization, and the proof template for almost-sure orthogonality.
RH-369 supplies the non-Parry branch-symmetric family and its parameter
identities.  RH-367 and the PCF parity factor are not used as a continuum
operator or as an arithmetic trace.  Gates A--E remain false/open.

\begingroup\sloppy
\begin{thebibliography}{9}
\raggedright
\bibitem{HenonMobius2026}
Source package, \emph{Möbius correlations on a certified Hénon survivor},
commit \texttt{34490443f50cfe9af9ff93888e51e7e7e534a5a7}, 2026.
\bibitem{RH366}
RH-366, \emph{Möbius orthogonality, adaptive encoding, and Parry covariance},
repository release, 2026.
\bibitem{Davenport1937}
H. Davenport, On some infinite series involving arithmetical functions (II),
\emph{Quart. J. Math.} 8 (1937), 313--320.
\bibitem{Tao2015}
T. Tao, The logarithmically averaged two-point Chowla and Elliott
conjectures, \emph{Forum Math. Pi} 3 (2015), e2.
\end{thebibliography}
\endgroup

\end{document}
